
\begin{center}
\renewcommand{\arraystretch}{1.5}
\begin{tabular}{ll}
\toprule
\textbf{Measure of center} & \textbf{Companion measure of spread} \\
\midrule
Mean   & Sample standard deviation \\
Median & Range \\
\bottomrule
\end{tabular}
\end{center}

\noindent
The mean and sample standard deviation are appropriate when the
data is approximately symmetric with no extreme values.
The median and range are more informative when the data
contains extreme values or is heavily skewed.

\begin{example}
The following data records the ages of participants in
two different programs.

Program A (yoga class):
\[ 28,\; 31,\; 34,\; 29,\; 33,\; 30,\; 32,\; 35. \]
\[
\bar{x}_A = 31.5, \qquad s_A \approx 2.45.
\]

Program B (community center drop-in):
\[ 8,\; 15,\; 22,\; 45,\; 67,\; 72,\; 78,\; 83. \]
\[
\bar{x}_B = 48.75, \qquad \text{Median}_B = 56, \qquad
\text{Range}_B = 75.
\]

Program A has a narrow age range and a small standard
deviation. The mean and sample standard deviation describe it well.
Program B spans a wide age range with no clear cluster.
The median and range are more informative for Program B.
\end{example}

\begin{exercises}
  \item For each data set, compute all five statistics:
    mean, median, mode, range, and sample standard deviation.
    Then state which pair (mean/SD or median/range) better
    describes the center and spread, and explain why.
    \begin{enumerate}[label=(\alph*), itemsep=6pt]
      \item $12,\; 14,\; 14,\; 15,\; 16,\; 17,\; 14$
      \item $5,\; 6,\; 7,\; 8,\; 9,\; 10,\; 85$
    \end{enumerate}

  \item A meteorologist records daily rainfall (in inches)
    for two weeks:
    \[ 0,\; 0,\; 0.2,\; 0,\; 1.4,\; 0,\; 0,\;
       0.1,\; 0,\; 0,\; 2.1,\; 0,\; 0,\; 0. \]
    \begin{enumerate}[label=(\alph*), itemsep=4pt]
      \item Find the mean, median, and mode.
      \item Which measure of center best represents a
        typical day?
        Justify in two sentences.
    \end{enumerate}

  \item A dataset has mean $20$ and sample standard deviation $3$.
    Another has mean $20$ and sample standard deviation $10$.
    \begin{enumerate}[label=(\alph*), itemsep=4pt]
      \item Sketch a rough histogram for each,
        placing the center at $20$.
      \item What does a larger sample standard deviation indicate
        about the distribution of values?
    \end{enumerate}
\end{exercises}

%--------------------------------------------------------------------
\section{Visualizing Center and Spread}
%--------------------------------------------------------------------

The histogram is the primary visual tool for understanding
the distribution of a data set.
Center and spread are visible in the shape of a histogram:
the center is where the distribution is tallest, and the
spread is how wide the distribution is.

\begin{example}
The following histogram shows the distribution of 50 exam
scores.

\begin{center}
\begin{tikzpicture}
\begin{axis}[
  ybar interval,
  width=10cm, height=5.5cm,
  bar width=1.0,
  xtick={50,60,70,80,90,100,110},
  xticklabels={50,60,70,80,90,100,{}},
  x tick label style={font=\small},
  xlabel={Score},
  ylabel={Frequency},
  y tick label style={font=\small},
  ytick={0,2,4,6,8,10,12,14},
  ymin=0, ymax=15,
  ymajorgrids=true,
  grid style=dotted,
]
\addplot[fill=mathblue!55, draw=mathblue]
  coordinates {(50,3)(60,7)(70,14)(80,13)(90,9)(100,4)(110,0)};
\end{axis}
\end{tikzpicture}
\end{center}

The distribution is approximately symmetric and centered
near $75$--$80$.
Most scores fall between $60$ and $90$.
A small number of scores appear below $60$ or above $90$.
\end{example}

\begin{example}[Two distributions with the same mean]
The two distributions below both have mean $50$,
but different sample standard deviations.

\begin{center}
\begin{tikzpicture}
\begin{axis}[
  width=10cm, height=6cm,
  xlabel={Value},
  ylabel={Relative Frequency},
  ytick={0,0.1,0.2,0.3,0.4},
  ymin=0, ymax=0.45,
  xmin=20, xmax=80,
  grid=major,
  grid style=dotted,
  legend style={at={(0.97,0.97)}, anchor=north east, font=\small},
]
\addplot[mathblue, very thick, smooth, samples=60,
  domain=22:78]
  {0.40*exp(-0.5*((x-50)/6)^2)};
\addlegendentry{Small $s$}
\addplot[mathred, very thick, smooth, dashed, samples=60,
  domain=22:78]
  {0.16*exp(-0.5*((x-50)/15)^2)};
\addlegendentry{Large $s$}
\draw[black, thick, dashed]
  (axis cs:50,0) -- (axis cs:50,0.40);
\node[above right, font=\small] at (axis cs:50,0.01)
  {mean $= 50$};
\end{axis}
\end{tikzpicture}
\end{center}

The narrow distribution has a small sample standard deviation:
values cluster tightly around the mean.
The wide distribution has a large sample standard deviation:
values are spread broadly around the mean.
Both are centered at the same location.
\end{example}

\begin{exercises}
  \item The histogram below shows the distribution of
    heights (in inches) for a sample of 40 adults.

  \begin{center}
  \begin{tikzpicture}
  \begin{axis}[
    ybar interval,
    width=9cm, height=5cm,
    bar width=1.0,
    xtick={58,62,66,70,74,78},
    xticklabels={58,62,66,70,74,{}},
    x tick label style={font=\small},
    xlabel={Height (inches)},
    ylabel={Frequency},
    y tick label style={font=\small},
    ytick={0,2,4,6,8,10,12},
    ymin=0, ymax=13,
    ymajorgrids=true,
    grid style=dotted,
  ]
  \addplot[fill=mathgreen!50, draw=mathgreen]
    coordinates {(58,2)(62,7)(66,12)(70,11)(74,6)(78,2)(82,0)};
  \end{axis}
  \end{tikzpicture}
  \end{center}

    \begin{enumerate}[label=(\alph*), itemsep=4pt]
      \item In which interval do most heights fall?
      \item Is the distribution symmetric or skewed?
      \item Would you expect the mean to be close to the
        median? Explain in one sentence.
    \end{enumerate}

  \item Sketch two histograms with the same center but
    different spreads. Label the center on both.
    Below each sketch, write one sentence describing
    what each distribution tells an observer about
    the data.

  \item A data set has mean $30$ and sample standard deviation $2$.
    A second data set has mean $30$ and sample standard deviation $8$.
    An observation of $36$ is taken.
    \begin{enumerate}[label=(\alph*), itemsep=4pt]
      \item How many sample standard deviations above the mean
        is $36$ in each data set?
      \item In which data set is the value $36$ more
        unusual? Explain in one sentence.
    \end{enumerate}
\end{exercises}

%--------------------------------------------------------------------
\section*{Chapter 3 Review}
%--------------------------------------------------------------------

\addcontentsline{toc}{section}{Chapter 3 Review}

\begin{enumerate}[label=\arabic*., leftmargin=2em, itemsep=8pt]

  \item The following data records the number of hours per
    week a sample of 10 students spend on homework:
    \[ 4,\; 7,\; 5,\; 8,\; 6,\; 5,\; 9,\; 4,\; 6,\; 6. \]
    \begin{enumerate}[label=(\alph*), itemsep=4pt]
      \item Find the mean, median, mode, and range.
      \item Construct the deviation table and compute
        the sample standard deviation.
      \item Describe the center and spread of this data
        set in two complete sentences.
    \end{enumerate}

  \item A company tracks the number of sick days taken
    by employees in one year.
    The data for eight employees is:
    \[ 0,\; 1,\; 2,\; 1,\; 0,\; 3,\; 1,\; 18. \]
    \begin{enumerate}[label=(\alph*), itemsep=4pt]
      \item Compute the mean and median.
      \item Which measure better represents a typical
        employee? Explain.
      \item If the value $18$ is removed, how do the mean
        and median each change? Compute both.
    \end{enumerate}

  \item Three students each score a mean of $80$ on five
    exams. Student A's scores are
    $78, 79, 80, 81, 82$.
    Student B's scores are
    $60, 70, 80, 90, 100$.
    Without computing, predict which student has the
    larger sample standard deviation.
    Then compute both sample standard deviations to verify.

  \item A data set has six values.
    Five of them are $12, 15, 14, 13, 16$.
    The mean of all six values is $14$.
    Find the sixth value.

  \item A histogram of test scores is skewed left.
    \begin{enumerate}[label=(\alph*), itemsep=4pt]
      \item What does it mean for a histogram to be
        skewed left?
      \item Would you expect the mean to be greater than,
        equal to, or less than the median?
        Explain in two sentences.
    \end{enumerate}

  \item The temperatures recorded over five days in two
    cities are:
    \begin{center}
    \renewcommand{\arraystretch}{1.4}
    \begin{tabular}{lrrrrr}
    \toprule
    & Day 1 & Day 2 & Day 3 & Day 4 & Day 5 \\
    \midrule
    City A & 68 & 70 & 69 & 71 & 67 \\
    City B & 55 & 62 & 74 & 79 & 85 \\
    \bottomrule
    \end{tabular}
    \end{center}
    \begin{enumerate}[label=(\alph*), itemsep=4pt]
      \item Compute the mean for each city.
      \item Compute the range and sample standard deviation
        for each city.
      \item Which city has more consistent temperatures?
        Justify using both spread measures.
    \end{enumerate}

  \item A quality control engineer samples 6 ball bearings
    and measures their diameters (in mm):
    \[ 9.98,\; 10.01,\; 10.00,\; 9.99,\; 10.02,\; 10.00. \]
    \begin{enumerate}[label=(\alph*), itemsep=4pt]
      \item Compute the mean and sample standard deviation.
      \item The specification requires that all bearings
        fall within $0.03$ mm of the target diameter of
        $10.00$ mm. Do all sampled bearings meet this
        requirement?
    \end{enumerate}

\end{enumerate}
